\documentclass[11pt,a4paper]{article}
\usepackage[margin=25mm]{geometry}
\usepackage[T1]{fontenc}
\usepackage{lmodern,microtype}
\usepackage{amsmath,amssymb,amsthm,mathtools}
\usepackage{booktabs,array,graphicx,xcolor}
\usepackage[colorlinks=true,linkcolor=blue!45!black,urlcolor=blue!55!black,citecolor=blue!45!black]{hyperref}
\usepackage{fancyhdr}
\hypersetup{pdftitle={Single-mode fidelity of collective photon emission},pdfauthor={Ruge Lin}}
\pagestyle{fancy}\fancyhf{}
\fancyhead[L]{\small Single-mode fidelity of collective emission}
\fancyhead[R]{\small Ruge Lin}
\fancyfoot[C]{\thepage}
\setlength{\headheight}{14pt}
\setlength{\parskip}{5pt}
\setlength{\parindent}{0pt}
\allowdisplaybreaks[1]
\raggedbottom
\newtheorem{theorem}{Theorem}
\newtheorem{proposition}[theorem]{Proposition}
\newtheorem{corollary}[theorem]{Corollary}
\newtheorem{lemma}[theorem]{Lemma}
\newcommand{\E}{\mathbb E}
\newcommand{\Prob}{\mathbb P}
\newcommand{\GE}{\Gamma_{\!E}}
\newcommand{\HH}{\mathsf H}
\newcommand{\KL}{D}
\newcommand{\norm}[1]{\left\lVert#1\right\rVert}
\newcommand{\Fock}[2]{\lvert #1_{#2}\rangle}
\newcommand{\dd}{\,\mathrm d}
\newcommand{\Tr}{\operatorname{Tr}}
\newcommand{\Bin}{\operatorname{Bin}}
\newcommand{\HG}{\operatorname{HG}}
\newcommand{\sech}{\operatorname{sech}}
\title{Single-mode fidelity of collective photon emission}
\author{Ruge Lin}
\date{17 October 2026}
\begin{document}
\maketitle
\begin{abstract}
For ideal Dicke emission from \(N\) initially excited emitters, the
maximum squared overlap with any single-temporal-mode \(N\)-photon
Fock state is \(\Theta(N^{-1/2})\). An explicit normalized pulse and
a relative-entropy inequality give the lower bound; temporal photon
counts give an upper bound uniform over all pulse shapes. The argument
extends to \(n\) symmetric excitations in \(M\) emitters: near full
inversion the fidelity is \(\Theta(\sqrt{(M-n+1)/M})\), whereas it
tends to one when \(n/M\to0\). Uniform attenuation retains a
conditional ceiling of order \(k^{-1/2}\) for \(k\) surviving photons,
and at fixed transmission the best unconditional photon-number and
mode target has yield \(\Theta(N^{-1})\). Independent local decay
instead produces a different exact conditional kernel and a
polynomial all-photon yield at fixed loss ratio. Diagonal collective
phases and common rate shaping cannot improve the ideal modal
fidelity. Detecting and consuming a prefix gives an exact resource
law, with explicit detection-efficiency and switching-delay costs.
Exact Fock emission from a coherent Markov ladder requires bosonic
rates and separable phases; under a native-rate cap exact purity is
slow. This exactness restriction is singular: a capped,
excitation-dependent coherent coupling achieves fidelity tending to
one at the native order of completion time, with polynomial
all-photon acceptance at fixed local loss. The control construction
specifies a mathematical coupling model rather than an experimental
implementation.
\end{abstract}

\section{Model, observables, and main result}
\label{sec:model}

We consider identical two-level emitters in their fully symmetric spin
sector, vacuum input, one desired Markov output, and complete eventual
emission. The native coupling is \(L_{\mathrm{out}}=\sqrt{\Gamma}J_-\),
with \(\Gamma>0\). Unless a section explicitly introduces loss or
control, there is no pumping, dephasing, feedback, or other output
channel. These are the ideal collective-emission assumptions used in
the waveguide descriptions of Refs.~\cite{paulisch,perarnau}.

Let \(b^\dagger(t)\) create an output photon, with
\([b(t),b^\dagger(s)]=\delta(t-s)\). For a normalized temporal mode,
\begin{equation}
 a_f^\dagger=\int_0^\infty f(t)b^\dagger(t)\dd t,\qquad
 \Fock{n}{f}=\frac{(a_f^\dagger)^n}{\sqrt{n!}}|0\rangle,\qquad
 F^\star(\rho_n)=\sup_{\norm{f}_2=1}\langle n_f|\rho_n|n_f\rangle.
 \label{eq:fidelitydef}
\end{equation}
Here \(\rho_n\) has trace one in the \(n\)-photon sector. All
fidelities are squared overlaps, or their mixed-state extension
against a pure target. The optimization includes arbitrary complex
pulse shapes and is not restricted to the most populated one-photon
mode. If the sector occurs with probability \(p_n\), its unconditional
target-state yield is \(Y_n^\star=p_nF^\star(\rho_n)\). Sector
probability, conditional mode fidelity, and yield are separate
quantities.

For the lossless native and partial-source calculations in
Sections~\ref{sec:wave}--\ref{sec:two}, time is measured in units
\(\Gamma^{-1}\), and \(t\) denotes the dimensionless coordinate.
Rescaling a physical mode by
\(\widetilde f(\tau)=\Gamma^{-1/2}f(\tau/\Gamma)\) preserves its norm
and all overlaps. Later loss and control sections use physical
\(t\), with dimensionless \(\tau=\Gamma t\) stated when needed.
Euler's gamma function is denoted \(\GE\), to distinguish it from
the emission rate.

\begin{theorem}[Fully inverted source]
For the ideal output \(\rho_N=|\Psi_N\rangle\langle\Psi_N|\) of \(N\) fully excited emitters, write \(F_N^\star=F^\star(\rho_N)\). Then
\begin{equation}\boxed{
 \sqrt{2\pi}\le\liminf_{N\to\infty}\sqrt N F_N^\star
 \le\limsup_{N\to\infty}\sqrt N F_N^\star
 \le\frac{4\sqrt{2\pi}}{e}.}
\label{eq:main}
\end{equation}
Consequently $F_N^\star=\Theta(N^{-1/2})$. The two endpoint constants are approximately $2.506628$ and $3.688548$.
\end{theorem}
Equation~\eqref{eq:main} determines the power, not the existence or value of a limiting prefactor. The displayed asymptotic constants are not finite-$N$ upper and lower inequalities without their limiting qualifications.

\section{Exact emitted wavefunction}
\label{sec:wave}
Use the normalized symmetric-coordinate convention
\begin{equation}
 |\Psi_N\rangle=\frac1{\sqrt{N!}}\int_{[0,\infty)^N}
 \psi_N(\boldsymbol t)\prod_{i=1}^N b^\dagger(t_i)\,d^N t,
 \qquad\int|\psi_N|^2\,d^Nt=1.
\end{equation}
The successive waiting times between emissions are independent exponentials with rates
\begin{equation}\lambda_j=j(N-j+1),\qquad j=1,\ldots,N.\end{equation}
For ordered times $0=t_0<t_1<\cdots<t_N$,
\begin{align}
 \sum_{j=1}^N\lambda_j(t_j-t_{j-1})
 &=\sum_{j=1}^N(2j-N)t_j\\
 &=\sum_i t_i+\sum_{i<j}|t_i-t_j|.
\end{align}
Since $\prod_j\lambda_j=(N!)^2$, the standard cascade amplitude~\cite{perarnau} becomes
\begin{equation}\boxed{
 \psi_N(\boldsymbol t)=\sqrt{N!}\exp\!\left[-\frac12\sum_i t_i-
 \frac12\sum_{i<j}|t_i-t_j|\right].}
\label{eq:wave}
\end{equation}
Indeed the ordered-domain density is $N!|\psi_N|^2=(N!)^2\exp[-\sum_j\lambda_j(t_j-t_{j-1})]$, which integrates to one over that domain. The symmetric density $P_N=|\psi_N|^2$ therefore integrates to one over the full orthant. Positivity permits replacing $f$ by $|f|$ without decreasing the overlap.

\section{Constructive lower bound}
\label{sec:lower}
\subsection{A classical affinity identity}
For a probability density $p$ on the half-line, put $f=\sqrt p$ and $Q_p(\boldsymbol t)=\prod_i p(t_i)$. The amplitude is exactly the classical Hellinger affinity,
\begin{equation}
 \langle N_f|\Psi_N\rangle=\int\sqrt{P_NQ_p}\,d^Nt
 =\E_{Q_p}\exp\!\left[\tfrac12\log(P_N/Q_p)\right].
\end{equation}
Jensen's inequality, followed by squaring, gives
\begin{equation}\boxed{F_N^\star\ge F_N[f]\ge\exp[-\KL(Q_p\Vert P_N)].}\end{equation}
From Eq.~\eqref{eq:wave},
\begin{equation}
 \KL(Q_p\Vert P_N)=N\int p\log p+N\E_p t+
 \frac{N(N-1)}2\E_p|t-s|-\log N!.
\label{eq:kl}
\end{equation}
No approximation or independence assumption about the actual output is used: only the trial density $Q_p$ is a product.

\subsection{An explicit normalized pulse}
Choose
\begin{equation}\boxed{
 p_N(t)=\frac{N^2e^{Nt}}{(e^{Nt}+N-1)^2},\qquad
 f_N(t)=\frac{Ne^{Nt/2}}{e^{Nt}+N-1},\qquad t\ge0.}
\label{eq:pulse}
\end{equation}
Its cumulative distribution is $(e^{Nt}-1)/(e^{Nt}+N-1)$, so normalization is exact. For $N>1$,
\begin{align}
 \E t&=\frac{\log N}{N-1},&
 \E|t-s|&=\frac{2[(N-1)-\log N]}{(N-1)^2},\\
 \int p_N\log p_N&=\frac{N\log N}{N-1}-2.
\end{align}
Substitution yields
\begin{equation}
 \KL(Q_{p_N}\Vert P_N)=\frac{N^2\log N}{N-1}-N-\log N!,
\end{equation}
and the exact finite-$N$ lower bound
\begin{equation}\boxed{
 F_N^\star\ge L_N:=\frac{e^N N!}{N^{N^2/(N-1)}}\quad(N\ge2).}
\label{eq:lower}
\end{equation}
For $N=1$, Eq.~\eqref{eq:pulse} reduces to $p_1(t)=e^{-t}$ and $F_1^\star=1$. Stirling's expansion gives
\begin{equation}
 L_N=\sqrt{\frac{2\pi}{N}}\exp\!\left[-\frac{\log N}{N-1}
 +\frac1{12N}+O(N^{-3})\right],
\end{equation}
which proves the lower half of Eq.~\eqref{eq:main}.

\subsection{Optimality within the relative-entropy variational principle}
The pulse in Eq.~\eqref{eq:pulse} globally minimizes Eq.~\eqref{eq:kl}, although it need not maximize the true fidelity. Let $X(u)$ be the quantile of $p$, $v(u)=X'(u)>0$, and $x_0=X(0)\ge0$. For an absolutely continuous quantile with finite objective,
\begin{equation}
 \frac{\KL(Q_p\Vert P_N)+\log N!}{N}
 =x_0+\int_0^1[A_N(u)v(u)-\log v(u)]\,du,
\end{equation}
where $A_N(u)=(1-u)[1+(N-1)u]$. The integrand is strictly convex in $v$ and minimized pointwise at $v=1/A_N$; $x_0$ is minimized at zero. Integration gives
\begin{equation}
 X_N(u)=\frac1N\log\frac{1+(N-1)u}{1-u},
\end{equation}
equivalent to Eq.~\eqref{eq:pulse}. Quantile jumps and singular increments add the nonnegative Stieltjes term $\int A_N\,dX_{\rm singular}$ and cannot improve the infimum. This establishes the same minimum for density-derived generalized quantiles with finite objective.

\section{Uniform upper bound from temporal photon counts}
\label{sec:upper}
\subsection{Every trial pulse supplies its own median cut}
Fix any normalized mode $f$. There is a time $T$ with $\int_0^T|f(t)|^2dt=1/2$. The number $m$ of photons before $T$ in $|N_f\rangle$ has distribution
\begin{equation}b_N(m)=2^{-N}\binom Nm.\end{equation}
Let $p_m(T)$ be the corresponding Dicke-output count distribution. Projection onto the orthogonal count sectors, followed by Cauchy--Schwarz within each sector, implies
\begin{equation}
 |\langle N_f|\Psi_N\rangle|
 \le\sum_{m=0}^N\sqrt{p_m(T)b_N(m)}.
\label{eq:countfidelity}
\end{equation}
The cut may depend on $f$ because the following estimates are uniform in $T$. This is an analysis device, not an assumption of feedback during emission.

\subsection{Arrival times bound count probabilities}
The emitted count is a pure-birth process with rates $r_m=(m+1)(N-m)$ from $m$ to $m+1$. This reduction is standard Dicke dynamics~\cite{malz}. Write $T_m$ for the $m$th arrival and $h_m$ for its density. For $1\le m\le N-1$,
\begin{equation}
 p_m(T)=\int_0^T h_m(s)e^{-r_m(T-s)}\,ds
 \le\frac{\norm{h_m}_\infty}{r_m}.
\label{eq:dwell}
\end{equation}
$T_m$ contains the independent first waiting time ${\rm Exp}(N)$, so convolution gives $\norm{h_m}_\infty\le N$. Already,
\begin{equation}
 p_m(T)\le\min\!\left\{1,\frac{N}{(m+1)(N-m)}\right\}.
\end{equation}
Together with Eq.~\eqref{eq:countfidelity}, this gives an explicit finite-$N$ upper bound by inserting the displayed expression for $1\le m<N$ and the bound one for $m=0,N$, and clipping the resulting squared sum at one.

To improve the asymptotic constant, fix $K\ge2$ independently of $N$. For $m\ge K$, convolution contraction gives $\norm{h_m}_\infty\le\norm{h_K}_\infty$. The density of $NT_K$ converges uniformly to that of $Y_K=\sum_{j=1}^K{\rm Exp}(j)$. This follows from convergence of each fixed-rate exponential density in both $L^1$ and $L^\infty$, followed by convolution. $Y_K$ has the distribution of the maximum of $K$ independent unit-rate exponentials, hence
\begin{equation}
 h_{Y_K}(s)=Ke^{-s}(1-e^{-s})^{K-1},\qquad
 \norm{h_{Y_K}}_\infty=c_K=(1-1/K)^{K-1}.
\end{equation}
Therefore $\norm{h_K}_\infty=N[c_K+o_N(1)]$. Uniformly over $T$ and $|m-N/2|\le N^{2/3}$,
\begin{equation}
 p_m(T)\le\frac{4c_K+o_N(1)}{N}.
\end{equation}

\subsection{The binomial width determines the exponent}
A binomial Chernoff bound and Cauchy--Schwarz give
\begin{equation}
 \sum_{|m-N/2|>N^{2/3}}\sqrt{p_m(T)b_N(m)}
 \le\sqrt2\,e^{-N^{1/3}}.
\end{equation}
Stirling's formula with Gaussian tail control gives
\begin{equation}
 \sum_{m=0}^N\sqrt{b_N(m)}=(2\pi N)^{1/4}[1+o(1)].
\end{equation}
Combining the central and tail estimates in Eq.~\eqref{eq:countfidelity}, uniformly over every $f$,
\begin{equation}
 F_N^\star\le\frac{4c_K\sqrt{2\pi}+o_N(1)}{\sqrt N}.
\end{equation}
First take $N\to\infty$ at each fixed $K$, then $K\to\infty$. Since $c_K\to e^{-1}$, the upper half of Eq.~\eqref{eq:main} follows. No exchange of a growing-$K$ limit is needed.


\subsection{Physical origin of the count bound}
The source starts at rate \(N\) and accelerates to rates of order
\(N^2\). The early waiting-time uncertainty is not erased by that
acceleration. At a trial mode's median cut, a Fock pulse has a
binomial count width of order \(\sqrt N\); Dicke emission has central
count probabilities at most of order \(1/N\). This mismatch bounds the
affinity even after pulse optimization. The arrival-density divided
by dwell-rate argument also applies to other radiative cascades,
although it alone does not supply a matching lower exponent for
another source.

\section{Dependence on the initial excitation}
\label{sec:partial}

Let \(M\) be the total number of emitters and \(n\) the excitation
number of an initially symmetric Dicke state. Put \(a=M-n+1\).
Exciting an arbitrary subset of \(n\) emitters in a larger array
does not in general prepare this state or preserve the symmetric
sector. In the lossless symmetric problem the emitted-count rates
are \(r_j=(n-j)(a+j)\), \(j=0,\ldots,n-1\). Their product is
\(n!\GE(M+1)/\GE(a)\), so the same ordered-coordinate normalization
as in Section~\ref{sec:wave} gives
\begin{equation}
 P_{M,n}(\boldsymbol t)=\frac{\GE(M+1)}{\GE(a)}
 \exp\!\left[-a\sum_i t_i-\sum_{i<j}|t_i-t_j|\right].
 \label{eq:affine}
\end{equation}
Its positive square root is the normalized emitted amplitude.

\subsection{An exact product relative-entropy bound}

The quantile argument of Section~\ref{sec:lower} now has
\(A(u)=(1-u)[a+(n-1)u]\), and its initial-offset term is \(a x_0\).
The pointwise minimizing derivative is again \(X'(u)=1/A(u)\).
For \(n>1\),
\begin{align}
 X_{M,n}(u)&=\frac1M
 \ln\frac{a+(n-1)u}{a(1-u)},\\
 p_{M,n}(t)&=\frac{aM^2e^{Mt}}{[ae^{Mt}+n-1]^2}.
 \label{eq:trial}
\end{align}
The cumulative distribution of this pulse is
\(a(e^{Mt}-1)/(ae^{Mt}+n-1)\). Thus normalization is exact.
Its product relative entropy with respect to \eqref{eq:affine} is
\begin{equation}
 \boxed{
 D_{M,n}=n\left[\frac{M\ln M-a\ln a}{n-1}-1\right]
 -\ln\GE(M+1)+\ln\GE(a),\qquad
 F_{M,n}^\star\ge e^{-D_{M,n}}.}
 \label{eq:D}
\end{equation}
The formula follows by substituting the quantile into the same
strictly convex functional used for the fully inverted source.
It is the global minimum of the product relative entropy, not
necessarily the pulse maximizing the actual fidelity. For \(n=1\),
the density is \(p(t)=M e^{-Mt}\), and fidelity one is exact.

Equation~\eqref{eq:affine}, its normalization, and the variational
calculation also hold for integer \(n\) and any real \(a\ge1\),
with \(M=n+a-1\). In that extension \(M\) is a distribution parameter,
not a fractional emitter number. This will be needed for local-loss
conditioning in Section~\ref{sec:local}.

\subsection{Uniform control in the dilute regime}

Set \(b=n-1\). Equation~\eqref{eq:D} is equivalently
\begin{equation}
 D_{M,n}
 =n\left[\frac1b\int_a^M\ln x\dd x
 -\frac1{b+1}\sum_{j=0}^b\ln(a+j)\right].
 \label{eq:gridentropy}
\end{equation}
The equally weighted, endpoint-inclusive grid dominates the uniform
interval in convex order. To see this, a convex function's integral
average is at most its composite trapezoidal average. If its grid
values are \(g_0,\ldots,g_b\), the grid average minus the trapezoidal
average is
\begin{equation}
 \frac{(b-1)(g_0+g_b)/2-\sum_{j=1}^{b-1}g_j}{b(b+1)}\ge0,
\end{equation}
where the inequality follows by bounding the internal values by
the endpoint chord. For \(x^2/2\), the grid average minus the
continuous average is \(b/12\). Comparing the curvature of
\(\ln x\) with \(-M^{-2}\) and \(-a^{-2}\), or applying convex
order to \(\ln x+x^2/(2a^2)\) and concave order to
\(\ln x+x^2/(2M^2)\), therefore gives
\begin{equation}
 \boxed{\frac{n(n-1)}{12M^2}\le D_{M,n}\le
 \frac{n(n-1)}{12a^2}.}
 \label{eq:dilutebound}
\end{equation}
Consequently,
\begin{equation}
 F_{M,n}^\star\ge
 \exp\!\left[-\frac{n(n-1)}{12(M-n+1)^2}\right],
 \qquad n/M\to0\ \Longrightarrow\ F_{M,n}^\star\to1.
 \label{eq:dilutefidelity}
\end{equation}
This remains uniform when \(n\to\infty\). The qualitative
weak-excitation route to Fock emission is established in
Ref.~\cite{liu}; \eqref{eq:dilutefidelity} quantifies the present
optimized-pulse problem.

\subsection{Fixed excitation fraction and uniform near-inversion scaling}

Write the Stirling remainder as
\begin{equation}
 \ln\GE(x)=(x-\tfrac12)\ln x-x+\tfrac12\ln(2\pi)+R(x),
 \qquad 0<R(x)<\frac1{12x},\quad x\ge1.
\end{equation}
Using \(\GE(M+1)=M\GE(M)\), Eq.~\eqref{eq:D} rearranges exactly to
\begin{equation}
 D_{M,n}=\frac12\ln\frac Ma-1+
 \frac a{n-1}\ln\frac Ma+R(a)-R(M).
 \label{eq:stirlingD}
\end{equation}
In particular, \(\ln(1+z)/z\le1\) gives the real-\(a\) bound
\begin{equation}
 F_{M,n}^\star\ge e^{-D_{M,n}}
 \ge e^{-1/(12a)}\sqrt{\frac aM}.
 \label{eq:realalower}
\end{equation}
If \(n/M\to\rho\in(0,1)\), Eq.~\eqref{eq:stirlingD} gives
\begin{equation}
 \liminf F_{M,n}^\star\ge
 e(1-\rho)^{(2-\rho)/(2\rho)}>0.
 \label{eq:fraction}
\end{equation}
This is a lower bound, not the exact limiting fidelity. At
\(\rho=1/2\) it is \(e/(2\sqrt2)=0.961057757\ldots\).

\begin{theorem}[Near complete inversion, including real \(a\)]
\label{thm:partial}
For integer \(n\), real \(a\ge1\), and \(M=n+a-1\), along every
sequence \(M\to\infty\) with \(a=o(M)\),
\begin{equation}
 \boxed{F_{M,n}^\star=\Theta\!\left(\sqrt{\frac aM}\right).}
 \label{eq:realascaling}
\end{equation}
The implicit constants may be chosen uniformly along these
sequences. For physical integer \(M,n\), this is the symmetric
partial-source fidelity.
\end{theorem}
\begin{proof}
The lower estimate is uniform by \eqref{eq:realalower}.
For the upper estimate take \(J=\max(2,\lceil a\rceil)\), which is
\(o(n)\). Each of the first \(J\) waiting rates is at most \(2na\).
The sum \(T_J\) of those waits therefore has characteristic function
bounded by
\begin{equation}
 |\widehat h_J(\omega)|
 \le[1+(\omega/(2na))^2]^{-J/2}.
\end{equation}
Fourier inversion gives
\begin{align}
 \norm{h_J}_\infty
 &\le\frac{na}{\sqrt\pi}
 \frac{\GE((J-1)/2)}{\GE(J/2)}
 \le\frac{na}{\sqrt{J-1}}
 \le\sqrt2\,n\sqrt a.
 \label{eq:fourierbound}
\end{align}
The middle inequality is the beta-integral bound
\(\int_{\mathbb R}(1+x^2)^{-J/2}\dd x\le\pi/\sqrt{J-1}\);
it follows from the Wallis gamma-ratio inequality and is equality
at \(J=2\). Convolution with later waits does not increase the
density supremum. The dwell bound \eqref{eq:dwell} now yields
\begin{equation}
 p_m(T)\le
 \frac{\sqrt2\,n\sqrt a}{(n-m)(a+m)},\qquad m\ge J.
\end{equation}
For \(n/4\le m\le3n/4\) this is \(O(\sqrt a/n)\).
At the median of an arbitrary trial pulse, the target count is
\(\Bin(n,1/2)\). Its square-root probabilities sum to
\(O(n^{1/4})\); the contribution outside the stated central
interval is exponentially small by Cauchy--Schwarz and the
binomial tail bound. Equation~\eqref{eq:countfidelity} therefore
gives \(F_{M,n}^\star=O(\sqrt{a/n})\).
Since \(n/M\to1\), this proves \eqref{eq:realascaling}.
\end{proof}

The trial lower bound has more explicit constants. If
\(a\to\infty\) while \(a=o(M)\), \eqref{eq:stirlingD} gives
\(e^{-D_{M,n}}=[e+o(1)]\sqrt{a/M}\). If \(a\) is fixed,
\begin{equation}
 \sqrt M\,e^{-D_{M,n}}\longrightarrow
 \sqrt{2\pi}\frac{a^a e^{1-a}}{\GE(a)}.
 \label{eq:fixedalower}
\end{equation}
These are asymptotics of the analytic lower bound, not asserted
optimal-fidelity prefactors. They include the fully inverted
case \(a=1\). At the opposite extreme \(n/a\to0\),
\eqref{eq:realalower} also implies fidelity tending to one.

\section{Exact optimal fidelity for two photons}
\label{sec:two}
For $N=2$, $\psi_2(t,s)=\sqrt2e^{-\max(t,s)}$ is a positive symmetric integral kernel. Its eigenvalue equation differentiates to
\begin{equation}f''+f'+\frac{\sqrt2}{\lambda}e^{-t}f=0,\qquad f'(0)=0.\end{equation}
Decay at infinity selects the Bessel solution. If $j_{0,1}$ is the first positive zero of $J_0$,
\begin{equation}\boxed{
 F_2^\star=\frac{32}{j_{0,1}^4}=0.956787875131927\ldots,\qquad
 f_2^\star(t)=\frac{e^{-t/2}J_1(j_{0,1}e^{-t/2})}{J_1(j_{0,1})}.}
\end{equation}
The pulse is exactly normalized. The first zero gives the largest kernel eigenvalue $4\sqrt2/j_{0,1}^2$.

\section{Uniform attenuation: surviving photons retain a modal ceiling}
\label{sec:attenuation}

Let \(\mathcal L_\eta\) be uniform, mode-independent attenuation of the
ideal output, with vacuum in the unused port and \(0<\eta<1\).
The number \(K_{\mathrm{out}}\) of transmitted photons is binomial:
\begin{equation}
 \Prob(K_{\mathrm{out}}=k)=B_{N,\eta}(k)
 =\binom Nk\eta^k(1-\eta)^{N-k}.
 \label{eq:binomialloss}
\end{equation}
Conditioned on \(k\), the temporal state is the normalized
\(k\)-particle reduced density operator \(\rho_N^{(k)}\). This follows
by applying the same beam splitter to each particle coordinate:
choosing the transmitted labels gives the binomial coefficient, and
tracing the other port traces the remaining coordinates. Write
\(F_{N\to k}^\star=F^\star(\rho_N^{(k)})\).

\begin{theorem}[Uniform survivor bound]
There is a universal \(C<\infty\) such that
\(F_{N\to k}^\star\le C/\sqrt{k}\) for all \(N\ge k\ge1\).
More precisely, along every sequence \(N\ge k\to\infty\),
\begin{equation}
 \limsup \sqrt{k}\,F_{N\to k}^\star
 \le\frac{4\sqrt{2\pi}}e.
 \label{eq:survivor}
\end{equation}
For a fixed retained fraction \(k/N\to\eta_0\in(0,1]\),
\(F_{N\to k}^\star=\Theta(N^{-1/2})\).
\end{theorem}

\begin{proof}
Let \(p_m(T)\) be the probability of \(m\) original photons before
a dimensionless cut \(T\). Section~\ref{sec:upper} already establishes
the arrival-density and dwell estimate
\begin{equation}
 p_m(T)\le\frac{N}{(m+1)(N-m)}
 \le\frac{N}{m(N-m)},\qquad 1\le m<N.
 \label{eq:entrydwell}
\end{equation}
The retained early count has the hypergeometric mixture
\begin{equation}
 q_j(T)=\sum_{m=0}^N p_m(T)
 \underbrace{\frac{\binom mj\binom{N-m}{k-j}}{\binom Nk}}
 _{\HG(N,m,k,j)}.
 \label{eq:hypergeom}
\end{equation}
For \(1\le j<k\), the kernel vanishes at \(m=0,N\). The identities
\(\binom mj/m=\binom{m-1}{j-1}/j\) and its complementary version,
followed by Vandermonde convolution, give
\begin{align}
 \sum_{m=1}^{N-1}\frac{\HG(N,m,k,j)}{m(N-m)}
 &=\frac{\binom{N-1}{k-1}}
 {j(k-j)\binom Nk}
 =\frac{k}{N j(k-j)},\\
 q_j(T)&\le\frac{k}{j(k-j)}.
 \label{eq:thinningbound}
\end{align}
This estimate is uniform in \(T,N,k\).

For any proposed mode choose a median time \(T\) of \(|f|^2\).
The target \(\Fock{k}{f}\) has cut-count law
\(b_j=\binom kj2^{-k}\). Monotonicity of fidelity under this
measurement implies
\begin{equation}
 \langle k_f|\rho_N^{(k)}|k_f\rangle
 \le\left(\sum_{j=0}^k\sqrt{q_j(T)b_j}\right)^2.
 \label{eq:measurement}
\end{equation}
On \(|j-k/2|\le k^{2/3}\), Eq.~\eqref{eq:thinningbound} is at most
\((4+o(1))/k\). Stirling's formula gives
\(\sum_j\sqrt{b_j}\sim(2\pi k)^{1/4}\). Outside this window, the
affinity contribution is at most the square root of the binomial
tail probability, by Cauchy--Schwarz, and is exponentially small in
\(k^{1/3}\). This proves a universal \(O(k^{-1/2})\) bound.

For the stated prefactor, Section~\ref{sec:upper} proves that
the density supremum of the first fixed \(J\ge2\) arrivals is
\(N[c_J+o_N(1)]\), where \(c_J=(1-1/J)^{J-1}\).
Thus the numerator \(N\) in \eqref{eq:entrydwell} can be replaced by
\(N[c_J+o_N(1)]\) whenever \(m\ge J\). For central retained
\(j\ge J\), every contributing \(m\) is at least \(J\), so the same
Vandermonde identity applies. Equation~\eqref{eq:measurement} yields
\(\limsup\sqrt{k}F_{N\to k}^\star\le4c_J\sqrt{2\pi}\).
Since \(N\ge k\to\infty\), the error is uniform along every stated
sequence. Taking \(J\to\infty\) gives \eqref{eq:survivor}.

For the lower bound, the falling-factorial occupation identity is
\begin{equation}
 \langle k_f|\rho_N^{(k)}|k_f\rangle
 =\left\langle\frac{(\hat n_f)_k}{(N)_k}\right\rangle
 \ge\langle N_f|\rho_N|N_f\rangle.
 \label{eq:rdmlower}
\end{equation}
The contribution from occupation \(n_f=N\) is one and all other
terms are nonnegative. Optimize and apply \eqref{eq:main}.
For \(k\) proportional to \(N\), the two orders match.
\end{proof}

\begin{corollary}[Optimize both output photon number and mode]
At fixed \(0<\eta<1\),
\begin{equation}
 \max_{0\le k\le N}\sup_f
 \langle k_f|\mathcal L_\eta(\rho_N)|k_f\rangle
 =\Theta(N^{-1}).
 \label{eq:anynumber}
\end{equation}
If all \(N\) photons are required instead, the yield is
\(\eta^NF_N^\star=\Theta(\eta^N/\sqrt N)\).
\end{corollary}
\begin{proof}
In a fixed bulk interval around \(\eta N\), the binomial point
probability is \(O(N^{-1/2})\) and the conditional fidelity is
\(O(N^{-1/2})\). Outside that interval the point probability is
exponentially small. For the reverse inequality choose a modal
binomial count, whose probability is \(\Theta(N^{-1/2})\), and use
\eqref{eq:rdmlower} with the full-source trial mode.
\end{proof}

A parasitic \emph{collective} channel \(\sqrt{\kappa}J_-\) is equivalent
to this beam-splitter loss after a clock rescaling, with
\(\eta=\Gamma/(\Gamma+\kappa)\). A time- or frequency-selective
filter is not uniform attenuation and is not covered by this theorem.

\section{Independent local decay and the cost of strong conditioning}
\label{sec:local}

Now add independent channels
\(L_i=\sqrt{\gamma}\,\sigma_i^-\), one for each emitter, with constant
\(\gamma\ge0\), and define \(\epsilon=\gamma/\Gamma\).
With no pumping, observing all \(N\) desired photons means that no
local jump occurred. The \(N\)-photon sector is pure: its local
environment is vacuum and the emitters end in their ground state.
It can be identified retrospectively by destructive total counting.
A usable herald that leaves the photons available requires an
additional nondestructive success flag or suitable monitoring of the
undesired environment.

\begin{proposition}[Exact all-photon sector]
Put \(a=1+\epsilon\) and \(M=N+a-1\). The sector probability is
\begin{equation}
 p_0=\prod_{j=1}^N\frac{j}{j+\epsilon}
 =\frac{\GE(N+1)\GE(1+\epsilon)}{\GE(N+1+\epsilon)}.
 \label{eq:localprob}
\end{equation}
Its unnormalized symmetric amplitude and normalized density are
\begin{align}
 \psi_{\mathrm{un}}(\boldsymbol\tau)
 &=\sqrt{N!}\,
   e^{-a\sum_i\tau_i/2-\sum_{i<j}|\tau_i-\tau_j|/2},\\
 P_{\mathrm{cond}}(\boldsymbol\tau)
 &=\frac{\GE(N+a)}{\GE(a)}
   e^{-a\sum_i\tau_i-\sum_{i<j}|\tau_i-\tau_j|}.
 \label{eq:localkernel}
\end{align}
The physical-time amplitude additionally has the Jacobian factor
\(\Gamma^{N/2}\).
\end{proposition}
\begin{proof}
Before desired jump \(j\), there are \(q=N-j+1\) excitations.
The desired rate is \(\Gamma jq\), and the sum of the local rates is
\(\gamma q\). The next jump is desired with probability
\(j/(j+\epsilon)\); its waiting time, independent of that jump
identity, has rate \(\Gamma q(j+\epsilon)\). Multiplying these
probabilities proves \eqref{eq:localprob}. Along a history with no
local jump, the integrated excitation number is \(\sum_i t_i\).
The no-local-jump propagator consequently multiplies the ideal
amplitude by \(e^{-\gamma\sum_i t_i/2}\), proving the first line.
Dividing its squared modulus by \eqref{eq:localprob} gives the
second line. Equivalently the normalized ordered waits have
dimensionless rates \((N-j+1)(a+j-1)\).
\end{proof}

The competing-rate, no-jump-product treatment of local decay is
used by Paulisch \emph{et al.} \cite{paulisch}. Its exact gamma
quotient supplies the success probability and normalization for
the optimized temporal-mode fidelity analyzed here.

\begin{corollary}[Conditional scaling and fixed-loss target yield]
For every sequence \(a=1+\epsilon\ge1\) with \(a=o(N)\),
\begin{equation}
 F_{\mathrm{cond}}^\star
 =\Theta\!\left(\sqrt{\frac aN}\right),
 \label{eq:localconditionalscaling}
\end{equation}
with uniform implicit constants. At fixed \(\epsilon\ge0\),
\begin{equation}
 p_0\sim\GE(1+\epsilon)N^{-\epsilon},\qquad
 Y_N^\star=p_0F_{\mathrm{cond}}^\star
 =\Theta(N^{-\epsilon-1/2}).
 \label{eq:fixedlossyield}
\end{equation}
\end{corollary}
\begin{proof}
The normalized kernel \eqref{eq:localkernel} is precisely
\eqref{eq:affine} with \(n=N\), \(a=1+\epsilon\), and
\(M=N+a-1\). The real-\(a\) theorem of Section~\ref{sec:partial}
therefore gives the conditional scaling, with
\(M/N\to1\) in this regime. The fixed-\(\epsilon\) gamma-ratio
asymptotic in \eqref{eq:localprob} gives the sector probability,
and multiplying proves the target-yield law.
\end{proof}

The unconditional optimized yield is nonincreasing in \(\epsilon\).
Indeed an optimal mode for the nonnegative kernel can be chosen
nonnegative, and increasing \(\epsilon\) decreases every such
unnormalized overlap pointwise. No corresponding monotonicity of
the normalized conditional fidelity is assumed.

\begin{corollary}[Strong conditioning has an exponential cost]
If \(\epsilon=cN\), with fixed \(c>0\), then
\begin{align}
 \liminf F_{\mathrm{cond}}^\star
 &\ge e\left(\frac c{1+c}\right)^{c+1/2}, \label{eq:strongF}\\
 p_0&\sim\sqrt{\frac{2\pi Nc}{1+c}}
 \left[\frac{c^c}{(1+c)^{1+c}}\right]^N. \label{eq:strongP}
\end{align}
More generally, keeping this native source's conditional fidelity
bounded below by a positive constant requires
\(\epsilon=\Omega(N)\), and hence exponentially small acceptance.
If \(\epsilon/N\to\infty\), the conditional fidelity tends to one.
\end{corollary}
\begin{proof}
For \(\epsilon=cN\), Eq.~\eqref{eq:stirlingD} tends to
\((c+\tfrac12)\ln[(1+c)/c]-1\); exponentiating proves
\eqref{eq:strongF}. Stirling's formula in \eqref{eq:localprob}
gives \eqref{eq:strongP}. If \(\epsilon\) were not \(\Omega(N)\),
a subsequence would have \(a/N\to0\), contradicting
\eqref{eq:realascaling}. When \(\epsilon\ge cN\), monotonicity of
\(p_0\) gives an exponential upper bound using \eqref{eq:strongP}.
Finally \(a/N\to\infty\) makes the lower bound
\eqref{eq:realalower} tend to one.
\end{proof}

At \(\epsilon=N\), these statements read
\(p_0=\binom{2N}{N}^{-1}\sim\sqrt{\pi N}\,4^{-N}\) and
\(\liminf F_{\mathrm{cond}}^\star\ge e/(2\sqrt2)
=0.961057757\ldots\).
For size-dependent \(\epsilon\), the exact gamma quotient must be
used outside the stated asymptotic regimes.

Local atomic loss is not uniform optical attenuation. After a local
jump the emitter state can leave the symmetric sector, and the
environment carries atom-specific information. For example, at
\(N=2\) the desired photon-number probabilities are
\begin{equation}
 \Prob(K_{\mathrm{out}}=2)=\frac2{(1+\epsilon)(2+\epsilon)},\quad
 \Prob(K_{\mathrm{out}}=1)=\frac{2\epsilon}{(1+\epsilon)(2+\epsilon)},
 \quad
 \Prob(K_{\mathrm{out}}=0)=\frac{\epsilon}{2+\epsilon}.
\end{equation}
To obtain them, after a first local jump decompose the remaining
one-excitation mixture equally into its symmetric bright and
antisymmetric dark states; only the bright part can emit into
\(J_-\). At \(\epsilon=1\) all three probabilities are \(1/3\),
which is not a two-photon binomial law. The survivor theorem of
Section~\ref{sec:attenuation} therefore does not apply to the
fewer-than-\(N\) sectors of local atomic loss.

The gamma-quotient formulas also extend to
\(\gamma(t)=\epsilon\Gamma(t)\) with constant ratio and complete
emission, by using \(\tau=\int_0^t\Gamma(s)\dd s\). If
\(\Gamma(t)\) is shaped while physical \(\gamma\) stays fixed,
the factor \(e^{-\gamma\sum_i t_i/2}\) remains in physical time;
the constant-ratio formulas cannot be substituted without a new
calculation.

\section{Controls that cannot improve the modal fidelity}
\label{sec:phases}

Label the symmetric ladder by the emitted count \(m=0,\ldots,N\).
Consider a single coherent output and a diagonal Hamiltonian,
\begin{equation}
 L(t)=\sum_{m=0}^{N-1}\sqrt{r_m(t)}e^{i\theta_m(t)}
 |m+1\rangle\langle m|,\qquad
 H(t)=\sum_{m=0}^{N}E_m(t)|m\rangle\langle m|.
 \label{eq:ladder}
\end{equation}
We set \(\hbar=1\). Every Hamiltonian that preserves both excitation
number and the fully symmetric spin sector has this diagonal
form, because each excitation sector of that spin representation
is one-dimensional. This statement does not cover Hamiltonians
that mix spin sectors or introduce a cavity, memory, or other
ancilla.

\begin{proposition}[Phase removal and scalar rate shaping]
At fixed rates \(r_m(t)\), removing the Hamiltonian and jump phases
cannot decrease the optimized full-photon fidelity or the optimized
fidelity of any normalized \(k\)-particle reduced state.
For native rates \(r_m(t)=\kappa(t)(m+1)(N-m)\), common deterministic
shaping with \(\int_0^\infty\kappa(t)\dd t=\infty\) leaves those
optimized fidelities unchanged.
\end{proposition}
\begin{proof}
Let \(\psi_0\ge0\) be the phase-free amplitude. Equation
\eqref{eq:ladder} changes it only to
\(\psi_H(\boldsymbol t)=\psi_0(\boldsymbol t)e^{i\Phi(\boldsymbol t)}\).
For retained variables \(\boldsymbol t\) and traced variables
\(\boldsymbol z\),
\begin{align}
 F_k(H;f)
 &=\int\dd\boldsymbol z\,
 \left|\int\psi_H(\boldsymbol t,\boldsymbol z)
                \prod_{i=1}^k f(t_i)^*\dd\boldsymbol t\right|^2\\
 &\le\int\dd\boldsymbol z\,
 \left[\int\psi_0(\boldsymbol t,\boldsymbol z)
                \prod_{i=1}^k |f(t_i)|\dd\boldsymbol t\right]^2
 =F_k(0;|f|).
 \label{eq:phaserm}
\end{align}
For \(k=N\) there are no traced variables. Optimization proves
the phase claim, including after uniform attenuation.

Let \(\tau(t)=\int_0^t\kappa(s)\dd s\). The map
\begin{equation}
 (U_\kappa f)(t)=\sqrt{\kappa(t)}\,f(\tau(t))
 \label{eq:isometry}
\end{equation}
is a one-photon isometry, and acts identically on every emitted
coordinate. The transformed field and all its reduced states lie
in its range, so optimization over all physical modes gives the
same fidelity as before the transformation. Partial tracing
commutes with this common isometry. This proves the rate claim.
\end{proof}

Arbitrary functions of \(J_z\), including nonlinear collective
phases within this ladder, are covered by \eqref{eq:phaserm}.
Combining them with common rate shaping cannot remove
\eqref{eq:main} or \eqref{eq:survivor}. The same invariance
holds for any common lossless one-photon unitary. Independent
local loss at a fixed physical rate changes the problem, as the
last paragraph of Section~\ref{sec:local} explains.

\subsection{Ideal all-photon mode filtering}
For an ideal filter retaining one normalized temporal mode \(f\)
and rejecting its orthogonal complement, the probability that
all \(N\) photons are retained is
\(|\langle N_f|\Psi_N\rangle|^2\). On that event the retained
state is \(\Fock{N}{f}\). Optimization gives
\(F_N^\star=\Theta(N^{-1/2})\) for the fully inverted source.
This equality follows by projecting onto occupation \(n_f=N\).
It gives an all-photon extraction probability, distinct from a
deterministic source guarantee. A usable herald must identify
success without consuming the desired retained photons; the
projection probability alone does not supply such an apparatus.

\section{Detect a prefix: an exact resource and imperfection law}
\label{sec:prefix}

Start with the lossless native source. Detect and consume its first
\(r\) photons, \(0\le r<N\), and immediately route all later photons
to the retained output. Use the \(r\)th-click time as the new time
origin; for \(r=0\), use the preparation time.

\begin{proposition}[Prefix identity]
With perfect detection, time resolution, and switching, the
record-aligned retained field has exactly the partial-source state
with \(M=N\), \(n=N-r\), \(a=r+1\). Hence
\begin{equation}
 F_{\mathrm{prefix}}(N,r)=F_{N,N-r}^\star.
 \label{eq:prefixidentity}
\end{equation}
In particular, for \(r=o(N)\),
\begin{equation}
 F_{\mathrm{prefix}}(N,r)
 =\Theta\!\left(\sqrt{\frac{r+1}{N}}\right).
 \label{eq:prefixlaw}
\end{equation}
If \(r/N\to\alpha\in(0,1)\), then
\begin{equation}
 \liminf F_{\mathrm{prefix}}(N,r)
 \ge e\,\alpha^{(1+\alpha)/[2(1-\alpha)]}.
 \label{eq:prefixfraction}
\end{equation}
\end{proposition}
\begin{proof}
Between clicks the symmetric ladder state acquires a scalar
no-jump factor; each click applies \(J_-\). After \(r\) desired
jumps the normalized state is the symmetric Dicke state with
\(N-r\) excitations, independent of all earlier waiting times.
The Markov property and stationary coupling make the future field
relative to that click precisely the partial-source output.
The remaining assertions follow from Section~\ref{sec:partial}. No favorable timing record is selected: in the ideal
model the \(r\)th click eventually occurs with probability one.
\end{proof}

A sublinear consumed prefix cannot keep the fidelity bounded away
from zero. A macroscopic prefix permits a nonzero limiting
fidelity while retaining a macroscopic photon number.
For \(\alpha=1/2\), the explicit lower bound is
\(e/(2\sqrt2)=0.961057757\ldots\), not the exact optimal fidelity.
If \(N-r=o(N)\), the dilute-source bound implies fidelity tending
to one. At finite size with \(N-r>1\), one may use
\(e^{-D_{N,N-r}}\). For one retained photon the fidelity is exactly
one. A simpler sufficient bound, valid also in that case, is
\begin{equation}
 F_{\mathrm{prefix}}(N,r)\ge
 \exp\!\left[-\frac{(N-r)(N-r-1)}{12(r+1)^2}\right].
\end{equation}
Ignoring the recorded origin produces a mixture of temporal
translations, a different target problem. This protocol consumes
the prefix; it is not a nondestructive measurement on the original
\(N\)-photon state, nor an optimality theorem for all adaptive controls.

\begin{proposition}[Exact nominal-target detection and latency cost]
Let \(r\ge1\), detector efficiency \(0<d\le1\), and a deterministic
delay \(\Delta\ge0\) between the \(r\)th detected click and routing.
Each prefix photon is independently detected, no other source
loss is present, and native emission continues during the delay.
For the predetermined target photon number \(n=N-r\), the optimized
record-aligned yield per initial trial is exactly
\begin{equation}
 Y_{N-r}^{\mathrm{nom}}(d,\Delta)
 =d^r e^{-\Gamma(N-r)(r+1)\Delta}F_{N,N-r}^\star.
 \label{eq:prefixcost}
\end{equation}
Conditioning the ensemble on a trigger divides this by
\(\Prob[\Bin(N,d)\ge r]\).
\end{proposition}
\begin{proof}
If \(S\) photons have actually been emitted at the \(r\)th detected
click, then
\begin{equation}
 \Prob(S=s,\mathrm{trigger})
 =\binom{s-1}{r-1}d^r(1-d)^{s-r},\qquad r\le s\le N.
\end{equation}
Summing gives the stated trigger probability. Only \(S=r\) can
leave the nominal \(N-r\) photons, and that event has probability
\(d^r\). From this post-click level, the next waiting rate is
\(\Gamma(N-r)(r+1)\). Keeping the same nominal number during the
delay requires no additional emission, with probability equal to
the exponential factor in \eqref{eq:prefixcost}. Conditioned on no
such jump, the normalized spin state is unchanged. Its future
field, referred to the switch time, is the same partial source.
Every other branch is orthogonal to the nominal photon-number
target, so multiplication gives the exact result.
\end{proof}

For \(r=\Theta(N)\), fixed \(d<1\) gives an exponential
nominal-yield penalty. To keep the detection factor nonvanishing
requires \(r[-\ln d]=O(1)\), equivalently \(r(1-d)=O(1)\) near
unit efficiency. When both \(r\) and \(N-r\) are proportional to
\(N\), the delay factor requires
\(\Gamma N^2\Delta=O(1)\); in general the exact condition is
\(\Gamma(N-r)(r+1)\Delta=O(1)\). The corresponding \(o(1)\) conditions
make the penalties tend to one. These statements apply to this
stopping rule and predetermined target. Extra number information,
accepting several output counts, buffering, or stopping emission
during the delay changes the protocol and its resource accounting.

\section{Exact single-mode ladders and an exact-purity time penalty}
\label{sec:exact}

We return to the coherent ladder \eqref{eq:ladder}, with locally
integrable rates, vacuum input, and complete emission. No extra
environment records its excitation-dependent control.
For a normalized target pulse \(p(t)=|f(t)|^2\), define its
survival \(S(t)=1-\int_0^t p(s)\dd s\) and hazard
\(h(t)=p(t)/S(t)\) where \(S(t)>0\).

\begin{theorem}[Exact Fock criterion within a Markov ladder]
The output is exactly \(\Fock{N}{f}\), up to a global phase, if and
only if on reachable times and levels
\begin{equation}
 r_m(t)=(N-m)h(t),\qquad
 \phi_m(t)=\arg f(t)+c_m\pmod{2\pi},
 \label{eq:exactcriterion}
\end{equation}
where \(c_m\) are constants and
\begin{equation}
 \phi_m(t)=\theta_m(t)
 -\int_0^t[E_m(s)-E_{m+1}(s)]\dd s .
 \label{eq:phase}
\end{equation}
The rate equality holds almost everywhere at every reachable time,
including pulse gaps where \(h=0\). The phase equality is required
only where \(f\ne0\).
When its cumulative probability is strictly between zero and one,
every level is reachable and the rate equality applies to every
\(m=0,\ldots,N-1\).
\end{theorem}
\begin{proof}
In an \(N\)-photon Fock mode the detection times are independent
draws from \(p\). Given \(m\) arrivals by time \(t\), the hazard of
the next is \((N-m)h(t)\). The physical ladder hazard is \(r_m(t)\).
Equality of the joint density therefore forces the first condition.
All levels have positive binomial probability when the cumulative
probability lies in \((0,1)\). Conversely, those hazards generate
the independent order-statistic law, so the positive amplitude
equals the target amplitude.

For ordered times \(0<t_1<\cdots<t_N\), the product of jump phases
and diagonal propagators has phase
\(\sum_{j=1}^N\phi_{j-1}(t_j)\), up to a global final-state phase.
To equal the target, its difference from
\(\sum_j\arg f(t_j)\) must be constant on the ordered support.
Varying one ordered coordinate at a time implies the second
condition almost everywhere; the constants are consistent across
overlapping ordered neighborhoods. The condition is sufficient
by multiplication. Support gaps and times after release impose
only the corresponding reachable-state requirements.
\end{proof}

For stationary positive rates and phases removed, the criterion
reduces to \(r_m=\kappa(N-m)\), with
\(f(t)=\sqrt{\kappa}e^{-\kappa t/2}\). A formal realization is
\begin{equation}
 L_{\mathrm{lin}}(t)
 =\sqrt{h(t)}\,J_-[N-\hat n+1]^{-1/2},
 \label{eq:linearjump}
\end{equation}
where \(\hat n\) is the current excitation number and the right
factor acts first. Its matrix element is \(\sqrt{h(t)n}\).
This is a truncated bosonic annihilation ladder. Bosonic
Fock-state release and pulse synthesis are established prior
constructions \cite{gough,brown,pfaff}; Eq.~\eqref{eq:linearjump}
identifies the coupling required inside this spin ladder.
Measuring the count and updating a scalar rate is not automatically
equivalent: an extra history record can destroy temporal coherence.

For reference, a stationary positive ladder has the following
exact overlap with the exponential comparison mode:
\begin{equation}
 F_{\exp}(\kappa)=
 \prod_{m=0}^{N-1}
 \left[\frac{2\sqrt{r_m\kappa(N-m)}}{r_m+\kappa(N-m)}\right]^2.
 \label{eq:stationaryoverlap}
\end{equation}
It follows by integrating the square root of each pair of
exponential waiting densities. Writing
\(r_m=\kappa(N-m)e^{\xi_m}\) gives
\(F_{\exp}=\prod_m\sech^2(\xi_m/2)\). Thus small independent
relative mismatches accumulate quadratically. This is a
comparison-mode formula, not a necessary error budget for the
fully optimized fidelity.

\begin{proposition}[Exact-purity penalty under a native-rate cap]
Assume the control can only attenuate native rates:
\begin{equation}
 0\le r_m(t)\le\Gamma(m+1)(N-m).
 \label{eq:ratecap}
\end{equation}
If its lossless output is an exact single-mode Fock state, then
\(h(t)\le\Gamma\) and
\begin{equation}
 \E T_{\mathrm{last}}\ge\frac{\HH_N}{\Gamma},
 \qquad \HH_N=\sum_{j=1}^N\frac1j.
 \label{eq:exactslow}
\end{equation}
Constant \(h=\Gamma\) attains this bound. Native emission has
\begin{equation}
 \E T_{\mathrm{last}}^{\mathrm{native}}
 =\frac{2\HH_N}{\Gamma(N+1)}.
 \label{eq:nativetime}
\end{equation}
If independent local loss \(\gamma\) is added to these same
lossless-exact controls, their all-desired sector remains
single-mode, but its probability is at most
\((1+\gamma/\Gamma)^{-N}\).
\end{proposition}
\begin{proof}
The \(m=0\) constraint in \eqref{eq:ratecap}, together with
\eqref{eq:exactcriterion}, gives \(h\le\Gamma\).
The zero-count state remains reachable whenever \(S(t)>0\), so
it cannot be removed from this argument by appealing to typical
counts. The cumulative hazard obeys
\(\int_0^t h\le\Gamma t\). Coupling each of the \(N\) independent
target arrival times to a unit-rate exponential shows that it is
no earlier than an exponential of rate \(\Gamma\). The expected
maximum of the latter is \(\HH_N/\Gamma\), proving
\eqref{eq:exactslow}. Summing the native waiting means
\(\{\Gamma j(N-j+1)\}^{-1}\) proves \eqref{eq:nativetime}.

With local loss, the unnormalized one-photon mode is
\begin{equation}
 v_\gamma(t)=e^{i\arg f(t)}\sqrt{h(t)}
 e^{-\frac12\int_0^t h(s)\dd s-\gamma t/2},\qquad
 q=\norm{v_\gamma}_2^2.
\end{equation}
The all-desired field is
\(q^{N/2}\Fock{N}{v_\gamma/\sqrt q}\), with probability \(q^N\).
Since its lossless arrival time is stochastically no earlier than
\(\operatorname{Exp}(\Gamma)\),
\(q=\E e^{-\gamma t}\le\Gamma/(\Gamma+\gamma)\).
Equality holds for constant \(h=\Gamma\).
\end{proof}

The exact lossless-purity constraint therefore incurs at least
a factor \((N+1)/2\) in mean completion time under the cap.
This statement does not classify controls designed only to purify
a lossy postselected sector, whose conditioning changes transition
hazards. It also does not extend to approximate purity. The next
construction shows why that distinction is essential.

\section{A capped coherent control with asymptotically unit fidelity}
\label{sec:capped}

Choose an even integer \(4\le K<N\) and set
\(T_0=-\Gamma^{-1}\ln(1-K/N)\). Define a target cumulative
distribution and its hazard by
\begin{align}
 u(t)&=
 \begin{cases}
  1-e^{-\Gamma t},&0\le t\le T_0,\\[2pt]
  \displaystyle\left[1+\frac{N-K}{K}
  e^{-\Gamma N(t-T_0)/2}\right]^{-1},&t\ge T_0,
 \end{cases} \label{eq:cappedcdf}\\
 h(t)&=\frac{u'(t)}{1-u(t)}
 =\begin{cases}\Gamma,&t<T_0,\\
 \Gamma Nu(t)/2,&t>T_0.
 \end{cases} \label{eq:cappedhazard}
\end{align}
The target mode is \(f=\sqrt{u'}\). Its ideal Fock ladder would
have rates \(s_m=(N-m)h\). Instead implement the single coherent
jump with zero phases and rates
\begin{equation}
 r_m^C(t)=(N-m)\min\{h(t),\Gamma(m+1)\}.
 \label{eq:cappedrates}
\end{equation}
They obey
\begin{equation}
 \Gamma(N-m)\le r_m^C(t)\le\Gamma(m+1)(N-m).
 \label{eq:cappedbounds}
\end{equation}
Every rung is attenuated from its native value by the amplitude
factor \(\min\{1,\sqrt{h(t)/[\Gamma(m+1)]}\}\). This is coherent,
excitation-dependent coupling to one output, not feedback based
on a measured photon record. The piecewise regular model permits
the hazard jump at \(T_0\).

\begin{theorem}[Finite-\(K\) fidelity and release-time guarantee]
Let
\begin{equation}
 \delta_K=\min\left\{1,\frac{e^{-K/8}}{1-e^{-1/4}}\right\}.
 \label{eq:deltak}
\end{equation}
The lossless source \eqref{eq:cappedrates} satisfies
\begin{equation}
 F_N^C[f]\ge(1-\delta_K)^2.
 \label{eq:cappedfidelity}
\end{equation}
For \(0<\delta_{\mathrm{tail}}<1\), define
\begin{equation}
 T_1=T_0+\frac2{\Gamma N}
 \ln\frac{N(N-K)}{K\delta_{\mathrm{tail}}}.
 \label{eq:T1}
\end{equation}
Then
\begin{equation}
 \E_C T_{\mathrm{last}}
 \le T_1+(\delta_K+\delta_{\mathrm{tail}})\frac{\HH_N}{\Gamma}.
 \label{eq:cappedtime}
\end{equation}
For the smallest even \(K\ge24\ln N\) and
\(\delta_{\mathrm{tail}}=N^{-3}\), at sufficiently large \(N\),
\begin{equation}
 F_N^C[f]\ge1-O(N^{-3}),\qquad
 \E_C T_{\mathrm{last}}
 =O\!\left(\frac{\ln N}{\Gamma N}\right).
 \label{eq:cappedasymptotic}
\end{equation}
\end{theorem}
\begin{proof}
Under the target Fock density, \(U_i=u(t_i)\) are independent
uniform variables. Write \(M(v)=\#\{i:U_i\le v\}\) and define
\begin{equation}
 G=\{M(v)+1\ge Nv/2\text{ for every }v\in[K/N,1]\}.
 \label{eq:good}
\end{equation}
On \(G\), actual and target rates agree at every time along the
entire history, including all no-jump intervals. Before \(T_0\)
they agree without restriction. Their positive amplitudes and
probability densities therefore coincide pointwise on \(G\).

If \(G\) fails at \(v\), let \(j=M(v)+1<Nv/2\).
When \(j\le K/2\), it follows that
\(M(K/N)<K/2\); the binomial lower-tail bound is \(e^{-K/8}\).
For \(j>K/2\), the \(j\)th order statistic satisfies
\(U_{(j)}>2j/N\), so a binomial with mean \(2j\) is at most
\(j-1\), of probability at most \(e^{-j/4}\).
Only \(j<N/2\) can violate the inequality, but extending the sum
to infinity increases the bound. Hence
\begin{equation}
 \Prob_f(G^c)
 \le e^{-K/8}+\sum_{j>K/2}e^{-j/4}
 =\frac{e^{-K/8}}{1-e^{-1/4}}\le\delta_K
\end{equation}
when the displayed untruncated bound is at most one; otherwise
the trivial probability bound supplies \(\delta_K=1\).
The normalized actual and target densities agree on \(G\), so
\(\Prob_C(G^c)=\Prob_f(G^c)\le\delta_K\).
Their affinity is at least their common probability on \(G\).
Squaring proves \eqref{eq:cappedfidelity}.

Equation~\eqref{eq:T1} gives
\(1-u(T_1)=\delta_{\mathrm{tail}}/(N+\delta_{\mathrm{tail}})\).
Thus the target probability of a late photon is
\(1-u(T_1)^N\le N[1-u(T_1)]\le\delta_{\mathrm{tail}}\).
Equality of the good-history densities implies
\(\Prob_C(T_{\mathrm{last}}>T_1)\le\delta_K+\delta_{\mathrm{tail}}\).
From any state and time, the lower rate bound in
\eqref{eq:cappedbounds} limits the mean remaining release time
by \(\HH_N/\Gamma\). Splitting the mean completion time at
\(T_1\) proves \eqref{eq:cappedtime}. Finally the stated choice
of \(K\) gives \(\delta_K=O(N^{-3})\), while
\(T_0=O(\ln N/(\Gamma N))\) and the second term in
\eqref{eq:T1} has the same order.
\end{proof}

The initial part of the target pulse contains mean count \(K\).
The later hazard rises with its deterministic cumulative
probability. Only strongly delayed count histories would need
a rate enhancement; clipping changes these rare histories.
Thus the exact-purity slowdown in Section~\ref{sec:exact} is not
a bound on fidelities tending to one. Equations
\eqref{eq:cappedrates}--\eqref{eq:cappedasymptotic} give a
constructive counterexample under the very same native-rate cap.

\subsection{Fixed local loss: polynomial acceptance and high conditional purity}

Add independent local decay at rate \(\gamma\) per excitation.
Let \(\psi_C\) be the normalized lossless capped-source amplitude.
The all-desired amplitude is
\(e^{-\gamma\sum_i t_i/2}\psi_C\), by the same integrated-excitation
identity used in Section~\ref{sec:local}. Define
\begin{equation}
 q(\gamma)=\int_0^\infty e^{-\gamma t}|f(t)|^2\dd t,\qquad
 g_\gamma(t)=q(\gamma)^{-1/2}e^{-\gamma t/2}f(t).
 \label{eq:tiltedmode}
\end{equation}

\begin{theorem}[Loss robustness of the capped construction]
Let \(P_{\mathrm{acc}}\) be the probability that all \(N\) photons
enter the desired output, and \(Y_N[g_\gamma]\) its unconditional
overlap with \(\Fock{N}{g_\gamma}\). Then
\begin{align}
 (1-\delta_K)q(\gamma)^N
 &\le P_{\mathrm{acc}}\le q(\gamma)^N, \label{eq:accbounds}\\
 Y_N[g_\gamma]&\ge q(\gamma)^N(1-\delta_K)^2,\\
 F_{N\mid\mathrm{all\ desired}}[g_\gamma]
 &\ge(1-\delta_K)^2. \label{eq:lossF}
\end{align}
Moreover, with \(\epsilon=\gamma/\Gamma\),
\begin{equation}
 q(\gamma)^N
 \ge\exp\{-\epsilon[K+2\ln(N/K)]\}.
 \label{eq:polysuccess}
\end{equation}
For fixed \(\epsilon\) and the \(K=O(\ln N)\) choice above,
acceptance has a polynomial lower bound in \(N\), while conditional
fidelity tends to one.
\end{theorem}
\begin{proof}
Exponential tilting by \(e^{-\gamma t}\) makes the distribution
\(|g_\gamma|^2\) stochastically earlier than \(|f|^2\). One may
see this by noting that an increasing indicator of a late time
has nonpositive covariance with the decreasing weight.
The event \(G\) is preserved by moving any arrival earlier.
An independent quantile coupling therefore gives
\(\Prob_{g_\gamma}(G)\ge1-\delta_K\).

On \(G\), \(\psi_C=\prod_i f(t_i)\). Integrating its locally
attenuated squared modulus over \(G\) gives
\(P_{\mathrm{acc}}\ge q^N\Prob_{g_\gamma}(G)\).
Likewise, restricting the positive overlap with
\(\prod_i g_\gamma(t_i)\) to \(G\) gives amplitude at least
\(q^{N/2}\Prob_{g_\gamma}(G)\), proving the yield bound.

For the upper bound on acceptance, \(r_m^C(t)\le s_m(t)\) at
every common state and time. Couple the two pure-birth processes
so that whenever their counts agree all actual jumps can be
shared with target jumps, with extra jumps allowed for the
target. When the counts differ, unit jumps cannot reverse their
order. This gives every actual arrival time no earlier than
its coupled target arrival. Hence
\begin{equation}
 P_{\mathrm{acc}}=\E_C e^{-\gamma\sum_i t_i}
 \le\E_f e^{-\gamma\sum_i t_i}=q^N .
\end{equation}
Dividing the yield lower bound by this upper bound proves
\eqref{eq:lossF}.

Integrating the survival function in \eqref{eq:cappedcdf} gives
the exact first moment
\begin{equation}
 N\E_f t
 =N\int_0^\infty[1-u(t)]\dd t
 =\frac{K+2\ln(N/K)}{\Gamma}.
 \label{eq:firstmoment}
\end{equation}
Jensen's inequality,
\(\E_f e^{-\gamma t}\ge e^{-\gamma\E_f t}\),
proves \eqref{eq:polysuccess}.
\end{proof}

The polynomial exponent in \eqref{eq:polysuccess} is a guarantee,
not an optimized loss frontier. It may be worse than the native
success exponent in \eqref{eq:fixedlossyield}. The result concerns
the all-desired sector; operational use of its conditioning has
the heralding qualification stated in Section~\ref{sec:local}.

\begin{lemma}[Accepted completion times are stochastically earlier]
For any such time-dependent pure-birth ladder with positive
cumulative rates and complete release, conditioning on no
independent local loss makes the completion time stochastically
no later than its lossless completion time. In particular,
\begin{equation}
 \E[T_{\mathrm{last}}\mid\mathrm{all\ desired}]
 \le\E_C T_{\mathrm{last}}.
 \label{eq:acceptedtime}
\end{equation}
\end{lemma}
\begin{proof}
Let \(R_m(t)=\int_0^t r_m(s)\dd s\). Using independent
unit-rate exponentials \(E_m\), construct successive arrivals by
\begin{equation}
 T_{m+1}=R_m^{-1}\bigl(R_m(T_m)+E_m\bigr),\qquad T_0=0.
\end{equation}
Here \(R_m^{-1}(x)=\inf\{t\ge0:R_m(t)\ge x\}\) is the
generalized inverse, allowing intervals with zero rate.
Each \(T_j\) is coordinatewise increasing in all the \(E_m\).
Independent product measures are associated: two increasing
functions have nonnegative covariance. This follows by the
one-variable Chebyshev covariance identity and induction over
independent coordinates. Therefore any increasing function of
the arrivals has nonpositive covariance with
\(W=e^{-\gamma\sum_jT_j}\), a decreasing function of those
coordinates. The no-local-loss conditional measure weights
histories by \(W/\E W\). Applying the covariance inequality to
\(\mathbf1_{\{T_N>t\}}\) proves stochastic domination; integrating
the tail inequality proves \eqref{eq:acceptedtime}.
\end{proof}

The argument uses monotonic cumulative rates and their inverses,
not monotonic instantaneous rates in time. Consequently accepted
runs of the capped construction also finish in mean time
\(O(\ln N/(\Gamma N))\).

\section{Reproducible numerical consistency checks}
\label{sec:numerics}

The executable \texttt{verification/check\_fidelity.py} evaluates
specified trial modes, verifies exact identities, and reports
ordinary floating-point convergence checks. Numerical overlaps
are achievable values for those modes, not certificates of
global optimality. The analytic uniform and asymptotic theorems
do not depend on fitting the numerical results.

For these numerical calculations set \(\Gamma=1\). For the native
source, use scaled time \(x=Nt\) and write
\(f(t)=\sqrt N\,\phi(x)\). Set \(c_0(0)=1\) and \(c_j(0)=0\)
for \(j>0\). The positive ordered-amplitude recursion is
\begin{equation}
 c_j'=-\frac{(j+1)(N-j)}{2N}c_j+
 j\sqrt{\frac{N-j+1}{N}}\,\phi(x)c_{j-1},
 \qquad j=0,\ldots,N,
 \label{eq:recursion}
\end{equation}
with the second term absent at \(j=0\). Then
\(c_N(\infty)\) is the exact trial amplitude. Thus an ordered
\(N\)-dimensional integral reduces to \(N+1\) ordinary differential
equations. The same construction with the general cascade rates
evaluates the partial-source and local-loss conditional kernels.
The reported values use the trial mode truncated at
\(x=X=\ln(M/a)+35\) and then exactly renormalized. If \(q\) is
the removed one-photon norm, then
\begin{equation}
 q=\frac{M e^{-X}}{a+(n-1)e^{-X}}.
\end{equation}
The infinite and truncated \(n\)-photon target states have squared
overlap \((1-q)^n\). Their projection probabilities against any
source therefore differ by at most
\(\sqrt{1-(1-q)^n}\), their pure-state trace distance. The code
reports this conservative analytic cutoff budget separately from
the measured change between two numerical cutoffs.

For the explicit fully inverted pulse \eqref{eq:pulse},
representative values are
\begin{center}
\begin{tabular}{@{}rrrr@{}}
\toprule
\(N\)&Trial fidelity&\(\sqrt N\) times trial fidelity&
Infinite-pulse lower bound\\
\midrule
1&1.000000000&1.000000&1.000000000\\
2&0.955818490&1.351731&0.923632012\\
5&0.833841887&1.864527&0.762238935\\
10&0.704689065&2.228422&0.618866167\\
20&0.565254318&2.527894&0.480738698\\
50&0.396241902&2.801853&0.327835432\\
100&0.292837776&2.928378&0.239469289\\
\bottomrule
\end{tabular}
\end{center}
The infinite one-photon pulse has fidelity one exactly. The exact
two-photon optimum in Section~\ref{sec:two} is larger than the
trial value, as required.

For the partial-source trial pulse \eqref{eq:trial}, the check
also evaluates the relative entropy by an independent
one-dimensional quantile integral:
\begin{center}
\begin{tabular}{@{}rrrr@{}}
\toprule
\(M\)&\(n\)&\(e^{-D_{M,n}}\)&Trial fidelity\\
\midrule
10&2&0.998151912&0.999058744\\
20&10&0.967523540&0.981861157\\
100&10&0.999176823&0.999581755\\
100&90&0.681685371&0.763737854\\
\bottomrule
\end{tabular}
\end{center}
A noninteger-parameter example is native emission with
\(N=5\), \(\epsilon=0.25\). Equation~\eqref{eq:localprob} gives
\(p_0=0.5883789413\); the conditional trial fidelity is
\(0.8693225961\), and the unconditional trial yield is
\(0.5114911087\).

Across these default rows, the largest change between the two
tolerance settings is \(2.7\times10^{-12}\), and the largest
change between tail parameters 30 and 35 is
\(6.5\times10^{-12}\). The largest mode-normalization residual is
\(2.2\times10^{-16}\), and the independent quantile-integral
discrepancy is \(5.7\times10^{-14}\). Evaluating the exact
two-photon Bessel mode with the cascade recursion gives a
squared-fidelity residual of \(1.9\times10^{-13}\).

The verification also compares the direct product and gamma
quotient for \(p_0\), the unweighted hypergeometric identity
\(\sum_m\HG(N,m,k,j)=(N+1)/(k+1)\), and the weighted identity
used in \eqref{eq:thinningbound}. The hypergeometric comparisons
use exact rational arithmetic. The program's machine-readable
output records the residuals, tolerance comparisons, and cutoff
budgets. These checks support implementation consistency; the
analytic proofs establish the uniform and asymptotic results.

\section{Relation to prior work and physical scope}
\label{sec:discussion}

The collective rate ladder, multiphoton wavefunction, and
multimode character of the emitted field are established
background. Paulisch and collaborators analyzed superradiant
photonic states and their metrological usefulness
\cite{paulisch}. Perarnau-Llobet and collaborators developed
efficient observables and explicit few-mode Fock decompositions
\cite{perarnau}. Lemberger and M\o lmer studied the one-photon
temporal modes and found a dominant occupation above 90 percent
\cite{lemberger}. That occupation fraction does not determine the
optimized probability that all photons share one mode.

Tziperman and collaborators studied complete propagating quantum
states and finite-size single-mode Fock fidelities, including
direct collective decay \cite{tziperman}. Liu and Shen's
weak-excitation Fock-generation regime \cite{liu} is consistent
with the dilute symmetric-source bounds above. The discrete
cavity-mode Tavis--Cummings setting of Ref.~\cite{yadav} is a
different output problem. Belliardo and collaborators optimize
temporal modes for information extraction from a superradiant
burst \cite{belliardo}; an information-optimal mode can contain
a vanishing fraction of the photons. A task-optimal mode therefore
need not be the most occupied one-photon mode.

Quantum-soliton analyses describe how center-of-mass
delocalization can suppress a many-body Hartree overlap even
when its \(N\)th root stays close to one \cite{holdaway}.
That root is distinct from a one-photon occupation or
reduced-state fidelity. The pair-distance factor in
Eq.~\eqref{eq:wave} makes this comparison relevant.
Malz, Trivedi, and Cirac established the large-\(N\) Dicke
population dynamics \cite{malz}; the count argument above
starts from that collective rate model and proves a uniform
arrival-density bound.

Deterministic optical Fock preparation, shaped bosonic input
synthesis, and controlled cavity release are established
constructions \cite{brown,gough,pfaff}. Dissipative Rydberg
filtering, photon subtraction, and temporal heralding provide
other state-preparation mechanisms \cite{gorshkov,pasharavesh,abudayyeh}.
They motivate distinctions among deterministic emission,
conditional fidelity, and unconditional target yield.

Transition-rate engineering along collective ladders is used,
for example, in the superabsorption proposal of Higgins and
collaborators \cite{higgins}. Resolving transition frequencies
and engineering their rates does not alone implement the
coherent single jump
\(\sum_m\sqrt{r_m(t)}|m+1\rangle\langle m|\) required by the
capped construction. Distinct output channels can retain rung
information. The construction specifies a coherent
excitation-dependent coupling to the same output, with no
additional environmental history record. It is not an
experimental architecture or a finite-bandwidth implementation.

The proved restrictions depend on the allowed operations.
Uniform attenuation, symmetric diagonal phases, and deterministic
common rate shaping retain the modal ceiling in their stated
senses. Prefix detection changes the photon-number resource.
Excitation-dependent coherent coupling changes the transition
ladder. Exact Fock factorization under a native-rate cap is slow,
while the controlled approximation in Section~\ref{sec:capped}
retains fast release and fidelity tending to one. These results
do not classify all excitation-preserving dynamics, other spin
sectors, ancillary memories, adaptive measurements, or non-Markov
sources. Emitter disorder, additional channels, and unobserved
local-loss histories require their own output-state analysis.

\begin{thebibliography}{99}\small
\bibitem{paulisch} V. Paulisch et al., \emph{Quantum metrology with one-dimensional superradiant photonic states}, Phys. Rev. A \textbf{99}, 043807 (2019). \href{https://arxiv.org/abs/1805.00712}{arXiv:1805.00712}.

\bibitem{perarnau} M. Perarnau-Llobet, A. Gonz\'alez-Tudela, and J. I. Cirac, \emph{Multimode Fock states with large photon number: effective descriptions and applications in quantum metrology}, Quantum Sci. Technol. \textbf{5}, 025003 (2020). \href{https://arxiv.org/abs/1910.03323}{arXiv:1910.03323}, especially Sec.~IV D and App.~I E.

\bibitem{lemberger} B. Lemberger and K. M\o lmer, \emph{Radiation eigenmodes of Dicke superradiance}, Phys. Rev. A \textbf{103}, 033713 (2021). \href{https://arxiv.org/abs/2009.11687}{arXiv:2009.11687}.

\bibitem{malz} D. Malz, R. Trivedi, and J. I. Cirac, \emph{Large-$N$ limit of Dicke superradiance}, Phys. Rev. A \textbf{106}, 013716 (2022). \href{https://arxiv.org/abs/2202.05197}{arXiv:2202.05197}.

\bibitem{tziperman} O. Tziperman et al., \emph{Nonlinear quantum light generation in collective spontaneous emission}, ACS Nano \textbf{19}, 21260--21270 (2025). \href{https://doi.org/10.1021/acsnano.4c15257}{doi:10.1021/acsnano.4c15257}; preprint \href{https://arxiv.org/abs/2306.11348}{arXiv:2306.11348} under its earlier title.

\bibitem{liu} Q. Liu and J.-T. Shen, \emph{Photonic Fock state generation using superradiance}, Opt. Lett. \textbf{47}, 4576--4579 (2022). \href{https://doi.org/10.1364/OL.468481}{doi:10.1364/OL.468481}.

\bibitem{holdaway} D. I. H. Holdaway, C. Weiss, and S. A. Gardiner, \emph{Quantum theory of bright matter-wave solitons in harmonic confinement}, Phys. Rev. A \textbf{85}, 053618 (2012). \href{https://arxiv.org/abs/1203.6006}{arXiv:1203.6006}, especially Sec.~V C.

\bibitem{yadav} A. Yadav and D. D. Yavuz, \emph{Quantum statistics of single-mode radiation emitted by superradiant Dicke states}, \href{https://arxiv.org/abs/2508.09962}{arXiv:2508.09962} (2025).

\bibitem{belliardo} F. Belliardo, A. Chu, M. Koppenh\"ofer, and A. A. Clerk, \emph{Extracting information from a superradiant burst using simple measurements}, \href{https://arxiv.org/abs/2603.13130}{arXiv:2603.13130v2} (31 March 2026).

\bibitem{gough}
J. E. Gough and G. Zhang, \emph{Generating nonclassical quantum
input field states with modulating filters}, EPJ Quantum Technol.
\textbf{2}, 15 (2015).
\href{https://doi.org/10.1140/epjqt/s40507-015-0027-z}
{doi:10.1140/epjqt/s40507-015-0027-z}.

\bibitem{brown}
K. R. Brown, K. M. Dani, D. M. Stamper-Kurn, and K. B. Whaley,
\emph{Deterministic optical Fock-state generation},
Phys. Rev. A \textbf{67}, 043818 (2003).
\href{https://doi.org/10.1103/PhysRevA.67.043818}
{doi:10.1103/PhysRevA.67.043818}.

\bibitem{pfaff}
W. Pfaff et al., \emph{Controlled release of multiphoton quantum
states from a microwave cavity memory}, Nature Phys. \textbf{13},
882--887 (2017).
\href{https://doi.org/10.1038/nphys4143}{doi:10.1038/nphys4143}.

\bibitem{pasharavesh}
A. Pasharavesh and M. Bajcsy, \emph{Multi-Photon Fock State
Generation via Selective Single Photon Subtraction in a Cascaded
Waveguide QED System}, Adv. Quantum Technol. \textbf{8}, 2400616
(2025). \href{https://doi.org/10.1002/qute.202400616}
{doi:10.1002/qute.202400616}.

\bibitem{gorshkov}
A. V. Gorshkov, R. Nath, and T. Pohl, \emph{Dissipative Many-Body
Quantum Optics in Rydberg Media}, Phys. Rev. Lett. \textbf{110},
153601 (2013).
\href{https://doi.org/10.1103/PhysRevLett.110.153601}
{doi:10.1103/PhysRevLett.110.153601}.

\bibitem{abudayyeh}
H. Abudayyeh, B. Lubotzky, S. Majumder, J. A. Hollingsworth, and
R. Rapaport, \emph{Purification of single photons by temporal
heralding of quantum dot sources},
\href{https://arxiv.org/abs/1805.00201}{arXiv:1805.00201} (2018).

\bibitem{higgins}
K. D. B. Higgins, S. C. Benjamin, T. M. Stace, G. J. Milburn,
B. W. Lovett, and E. M. Gauger, \emph{Superabsorption of light via
quantum engineering}, Nature Commun. \textbf{5}, 4705 (2014).
\href{https://doi.org/10.1038/ncomms5705}{doi:10.1038/ncomms5705}.
\end{thebibliography}
\end{document}
